\section{Recitation 6: Mutable State and Tables}
\label{sec:rec6}

This section reviews the store-passing design, makes evaluation order visible, and then adds a new kind of value, an immutable \emph{table} (key-value map), to see which things a store makes changeable and which it does not. It ends with a rewrite of part of the interpreter in which the store-passing is done by two helper functions.

\subsection{The Language}

The language has numbers, \code{+}, \code{*}, \code{let}, \code{lambda}, application, boxes, \code{begin}, and mutable variables with \code{set!}, plus four forms for tables:
\begin{lstlisting}[style=repl]
(make-table (k1 e1) (k2 e2) ...)   a new table mapping each key ki to the value of ei
(get-table t k)                    the value for key k in table t
(set-table t k v)                  a NEW table: like t, but with k mapped to v
(remove-table t k)                 a NEW table: like t, but without k
\end{lstlisting}
Two design decisions are fixed:
\begin{itemize}
  \item \textbf{Keys are symbols written in the program}, like field names, not expressions. \code{(get-table t (+ 1 2))} is not allowed; the key is known when the program is parsed.
  \item \textbf{Tables never change.} \code{set-table} and \code{remove-table} leave \code{t} alone and return a new table, like Racket's \code{hash-set} and \code{hash-remove}. Most of \Cref{sec:rec6-vs} follows from this.
\end{itemize}
\begin{lstlisting}
(define-type Entry
  [entry (k : symbol) (e : Expr)])

(define-type Expr
  [numC (n : number)]
  [plusC (e1 : Expr) (e2 : Expr)]
  [timesC (e1 : Expr) (e2 : Expr)]
  [letC (x : symbol) (e1 : Expr) (e2 : Expr)]
  [lambdaC (x : symbol) (e : Expr)]
  [appC (e1 : Expr) (e2 : Expr)]
  [idC (x : symbol)]
  [boxC (e : Expr)]
  [unboxC (e : Expr)]
  [setboxC (e1 : Expr) (e2 : Expr)]
  [setC (x : symbol) (e2 : Expr)]
  [seqC (e1 : Expr) (e2 : Expr)]
  [maketableC (entries : (listof Entry))]
  [gettableC (tbl : Expr) (k : symbol)]
  [settableC (tbl : Expr) (k : symbol) (val : Expr)]
  [removetableC (tbl : Expr) (k : symbol)])

(define-type-alias Table (hashof symbol Value))

(define-type Value
  [numV (n : number)]
  [closV (env : Env) (x : symbol) (e : Expr)]
  [boxV (l : Location)]
  [tableV (t : Table)])
\end{lstlisting}
Compare the two new kinds of value. A \code{boxV} holds a \emph{location}: the contents live in the store, so they can change while the \code{boxV} itself stays the same. A \code{tableV} holds the whole \tyname{Table} directly. There is nothing in it that points into the store, so a table value, once made, is fixed.

\tyname{Entry} is a separate type for the same reason as \tyname{PatternE} in \Cref{sec:rec4}: one \code{(k e)} pair carries a symbol and an expression, and a \code{maketableC} has a list of them.
\begin{lstlisting}
(define (parse-entry (s : s-expression)) : Entry
  (let ([l (s-exp->list s)])
    (entry (s-exp->symbol (first l)) (parse (second l)))))
\end{lstlisting}
\code{parse} gets one line per new form:
\begin{lstlisting}
            [(set!) (setC (s-exp->symbol (second l)) (parse (third l)))]
            [(make-table) (maketableC (map parse-entry (rest l)))]
            [(get-table) (gettableC (parse (second l)) (s-exp->symbol (third l)))]
            [(set-table) (settableC (parse (second l)) (s-exp->symbol (third l))
                                    (parse (fourth l)))]
            [(remove-table) (removetableC (parse (second l)) (s-exp->symbol (third l)))]
\end{lstlisting}
The keys go through \code{s-exp->symbol}, not \code{parse}: that is where ``keys are not expressions'' is enforced. The two \code{else} branches of \code{parse} send everything that is not a keyword to application, whether the head is a variable like \code{(f 10)} or an expression like \code{((get-table t inc) 41)}.

\subsection{Environments and the Store}

With \code{set!}, the meaning of a name is split over two maps:
\begin{lstlisting}[style=repl]
env   : symbol   -> Location      (fixed when the binding is made)
store : Location -> Value         (can be overridden later)
\end{lstlisting}
The store, \code{fetch}, \tyname{Result}, and \code{new-loc} are those of \Cref{sec:storepass}. What changes is the environment: a binding now holds a location, and \code{lookup} returns it.
\begin{lstlisting}
(define-type Binding
  [bind (name : symbol) (loc : Location)])
(define-type-alias Env (listof Binding))

(define (lookup (x : symbol) (env : Env)) : Location
  (cond
    [(cons? env)
     (if (equal? (bind-name (first env)) x)
         (bind-loc (first env))
         (lookup x (rest env)))]
    [else (error 'lookup "No binding found")]))
\end{lstlisting}
In \code{eval-env}, the cases for numbers, \code{lambda}, \code{+}, \code{*}, \code{let}, \code{begin}, \code{box}, and \code{unbox} are as in \Cref{sec:storepass,sec:storeboxes}. These four differ:
\begin{lstlisting}
    [appC (e1 e2)
          (type-case Result (eval-env env sto e1)
            [res (v1 sto-1)
                 (type-case Result (eval-env env sto-1 e2)
                   [res (v2 sto-2)
                        (let ([l (new-loc)])
                          (eval-env
                           (extend-env (bind (closV-x v1) l) (closV-env v1))
                           (override-store (cell l v2) sto-2)
                           (closV-e v1)))])])]
    [setC (x e)
          (type-case Result (eval-env env sto e)
            [res (v sto-1)
                 (let [(l (lookup x env))]
                   (res v (override-store (cell l v) sto-1)))])]
    [idC (x) (res (fetch (lookup x env) sto) sto)]
    [setboxC (e1 e2)
             (type-case Result (eval-env env sto e1)
               [res (v1 sto-1)
                    (type-case Result (eval-env env sto-1 e2)
                      [res (v2 sto-2)
                           (res v2 (override-store (cell (boxV-l v1) v2) sto-2))])])]
\end{lstlisting}
What to notice:
\begin{itemize}
  \item \textbf{Every application allocates.} \code{appC} gets a fresh location for the parameter and stores the argument there, and \code{letC} is application in disguise. So every variable has its own cell, and that is what lets \code{set!} work.
  \item \textbf{\code{set!} and \code{set-box!} both override a cell}, but they find it differently. \code{setC} looks the \emph{name} up in the environment; \code{setboxC} takes the location out of the \emph{value} it evaluated. \code{set!} can only reach the variable it names; a box can be passed around, and whoever holds it can change it.
  \item \textbf{\code{set-box!} returns the new contents} \code{v2} here, where the interpreter of \Cref{sec:storeboxes} returns the box \code{v1}. Of the examples below, only \code{ex2} sees the difference.
  \item \textbf{The store is threaded, never shared.} Each case passes the store that the previous subexpression returned to the next one. Which subexpression goes first is a choice the interpreter makes (\Cref{sec:rec6-order}).
\end{itemize}
\begin{lstlisting}
(define (eval (e : Expr)) : Result
  (eval-env empty-env empty-store e))

(define (run (s : s-expression)) : Value
  (res-v (eval (parse s))))
\end{lstlisting}

\subsection{Warm-up}

\begin{exercise}[The examples of \Cref{sec:stateboxes}, and \code{set!}]
Predict the value of each program before running it.
\begin{lstlisting}
(define ex0 '(let v (+ 1 1) (+ v v)))
(define ex1 '(begin (+ 1 1) (+ 2 2)))
(define ex2 '(let b (box 0) (set-box! b 1)))
(define ex3 '(let b (box 0) (unbox b)))
(define ex4 '(let b (box 0) (begin (begin (set-box! b 1)
                                         (set-box! b (+ 1 (unbox b))))
                                  (unbox b))))
(define ex5 '(let a (box 1)
               (let f (lambda x (+ x (unbox a)))
                 (begin
                   (set-box! a 2)
                   (f 10)))))
(define ex6 '(let v (+ 1 1) (begin (set! v 3) (+ v v))))
\end{lstlisting}
\end{exercise}
\begin{solution}
\begin{lstlisting}
(test (run ex0) (numV 4))
(test (run ex1) (numV 4))
(test (run ex2) (numV 1))
(test (run ex3) (numV 0))
(test (run ex4) (numV 2))
(test (run ex5) (numV 12))
(test (run ex6) (numV 6))
\end{lstlisting}
\begin{center}
\begin{tabular}{@{}lll@{}}
\toprule
ex & result & why \\
\midrule
ex1 & 4 & \code{begin} throws away the first value \\
ex2 & 1 & \code{set-box!} returns the new contents (\code{v2}) \\
ex4 & 2 & the second \code{set-box!} reads the store the first one returned \\
ex5 & 12 & \code{f} reads the box when it is \emph{called}: $10 + 2$ \\
ex6 & 6 & \code{set!} overrides \code{v}'s cell, then $3 + 3$ \\
\bottomrule
\end{tabular}
\end{center}
\code{ex5} is the one to understand. The closure for \code{f} captures the environment, which maps \code{a} to a location. The closure does not capture the store. By the time \code{(f 10)} runs, the store says that \code{a}'s box holds 2.
\end{solution}

\subsection{Evaluation Order}
\label{sec:rec6-order}

Our interpreter evaluates \code{e1} before \code{e2} in \code{(+ e1 e2)}, because it has to pick one to receive the store first. With side effects, the choice is visible:
\begin{lstlisting}
(define ex-order
  '(let b (box 10)
     (+ (begin (set-box! b 5) 4)
        (unbox b))))

(test (run ex-order) (numV 9))
\end{lstlisting}
\begin{lstlisting}[style=repl]
left to right:  (begin (set-box! b 5) 4) runs first, so (unbox b) sees 5:   4 + 5  = 9
right to left:  (unbox b) runs first and sees 10:                           4 + 10 = 14
\end{lstlisting}
Real languages make different choices. The same program can be written in C, Racket, and OCaml: \code{first()} adds 1 to a global counter, \code{second()} adds 10, and each prints the counter. Both are called once as the arguments of a function and once as the operands of \code{+}. The C version:
\begin{lstlisting}[language=C]
#include <stdio.h>

int count = 0;

__attribute__((noinline))
int first(void) {
    count += 1;
    printf("[first]  count = %d\n", count);
    return 10;
}

__attribute__((noinline))
int second(void) {
    count += 10;
    printf("[second] count = %d\n", count);
    return 20;
}

__attribute__((noinline))
void print_sum(int a, int b) {
    printf("Result of a + b: %d\n", a + b);
}

int main(void) {
    printf("Demonstration 1: Function Argument Evaluation Order\n");
    print_sum(first(), second());

    count = 0; // reset count

    printf("Demonstration 2: Expression Operand Evaluation Order\n");
    int total = first() + second();
    printf("Total: %d\n", total);

    return 0;
}
\end{lstlisting}
Compiled with Apple clang 21, at \code{-O0} and \code{-O2}, it prints:
\begin{lstlisting}[style=repl]
Demonstration 1: Function Argument Evaluation Order
[first]  count = 1
[second] count = 11
Result of a + b: 30
Demonstration 2: Expression Operand Evaluation Order
[first]  count = 1
[second] count = 11
Total: 30
\end{lstlisting}
\begin{center}
\begin{tabular}{@{}lll@{}}
\toprule
language & order & guaranteed? \\
\midrule
Racket & left to right & yes, by the language definition \\
C & Apple clang 21, \code{-O0} and \code{-O2}: left to right & no: unspecified \\
OCaml & right to left in practice & no: unspecified \\
\bottomrule
\end{tabular}
\end{center}
In C the order is ``arbitrary'': the standard leaves it unspecified for both arguments and operands, so another compiler (GCC on x86 is the usual example) may print \code{[second]} first. The sum is 30 either way. Only the trace, and anything else the calls do to \code{count}, depends on the order. (The Racket version never resets \code{count}, so its second demonstration prints 12 and 22 instead of 1 and 11; the order is the same.) The OCaml version is named for OCaml's usual right-to-left order; it was not run here.

The lesson for interpreters: the language designer \emph{chooses} the order, and in a store-passing interpreter the choice is written down explicitly by which \code{eval-env} call gets \code{sto} and which gets \code{sto-1}. In \Cref{sec:rec6-vs}, \code{ex-f} checks the order for \code{set-table}.

\subsection{Table Helpers}

A \tyname{Table} is a Racket hash. The four operations we need:
\begin{lstlisting}[style=repl]
(hash empty)       an empty hash
(hash-ref t k)     (some v) if k maps to v in t, (none) otherwise
(hash-set t k v)   a new hash: like t, but with k mapped to v
(hash-remove t k)  a new hash: like t, but without k
\end{lstlisting}

\begin{exercise}[\code{expect-table}]
Write \code{(expect-table v)}, which returns the table inside \code{v}, or signals the error \code{"expected a table"} if \code{v} is not a table. Hint: \code{type-case} can have an \code{else} case, for ``every other kind of value''.
\end{exercise}
\begin{solution}
\begin{lstlisting}
(define (expect-table (v : Value)) : Table
  (type-case Value v
    [tableV (t) t]
    [else (error 'expect-table "expected a table")]))
\end{lstlisting}
\code{eval-env} gets back a \tyname{Value}, and only a \code{tableV} can be used as a table. Instead of repeating that check in three cases, it gets its own helper. Without it we would write \code{(tableV-t v)}, and a non-table would fail with a message about the accessor instead of about the program.
\end{solution}

\begin{exercise}[\code{make-table}]
Write \code{(make-table env sto entries acc)}, which evaluates each entry's expression, left to right, passing the store along, and adds each key and its value to \code{acc}. It returns the finished table, along with the final store. If a key appears twice, the later entry wins. Hint: look at how the \code{plusC} case passes the store from \code{e1} to \code{e2}.
\end{exercise}
\begin{solution}
This is the only helper that runs code. It is the \code{plusC} pattern stretched over a list, plus an accumulator for the table built so far:
\begin{lstlisting}
(define (make-table (env : Env) (sto : Store) (entries : (listof Entry)) (acc : Table))
  : Result
  (cond
    [(empty? entries) (res (tableV acc) sto)]
    [else
     (type-case Entry (first entries)
       [entry (k e)
              (type-case Result (eval-env env sto e)
                [res (v sto-1)
                     (make-table env sto-1 (rest entries) (hash-set acc k v))])])]))
\end{lstlisting}
\begin{itemize}
  \item The recursive call gets \code{sto-1}, not \code{sto}. Passing \code{sto} would quietly drop the first entry's side effects; the \code{set-box!} test in \Cref{sec:rec6-interp} catches that.
  \item ``The later entry wins'' for a repeated key comes for free: \code{hash-set} on an existing key replaces it, and later entries are added later.
  \item It returns a \tyname{Result}, not a \tyname{Table}, because the caller needs the final store too.
\end{itemize}
\end{solution}

\begin{exercise}[\code{get-table}, \code{set-table}, \code{remove-table}]
Write \code{(get-table t k)}, which returns the value for \code{k} in \code{t} or signals \code{(string-append "no such key: " (symbol->string k))}; \code{(set-table t k v)}, a new table with \code{k} mapped to \code{v}; and \code{(remove-table t k)}, a new table without \code{k}, which signals the same error as \code{get-table} if \code{t} has no key \code{k}. Hint: which helper already signals that error?
\end{exercise}
\begin{solution}
These three are pure: no store, no evaluation, just hash operations. \code{hash-ref} returns an \code{optionof}, taken apart with \code{type-case} as in \Cref{sec:rec5}:
\begin{lstlisting}
(define (get-table (t : Table) (k : symbol)) : Value
  (type-case (optionof Value) (hash-ref t k)
    [some (v) v]
    [none () (error 'get-table
                    (string-append "no such key: " (symbol->string k)))]))

(define (set-table (t : Table) (k : symbol) (v : Value)) : Table
  (hash-set t k v))

(define (remove-table (t : Table) (k : symbol)) : Table
  (begin (get-table t k)
         (hash-remove t k)))
\end{lstlisting}
\code{hash-remove} on a missing key silently returns the same hash, but \code{remove-table} must raise an error. \code{get-table} already signals it, so \code{remove-table} calls it only for its error and throws the value away with \code{begin}.
\begin{lstlisting}
(test (expect-table (tableV (hash empty))) (hash empty))
(test/exn (expect-table (numV 1)) "expected a table")

(define t-ab (hash (list (values 'a (numV 1)) (values 'b (numV 2)))))
(test (get-table t-ab 'b) (numV 2))
(test/exn (get-table t-ab 'c) "no such key: c")
(test (get-table (set-table t-ab 'c (numV 3)) 'c) (numV 3))
(test (get-table (set-table t-ab 'a (numV 3)) 'a) (numV 3))
(test (get-table t-ab 'a) (numV 1))
(test/exn (get-table (remove-table t-ab 'a) 'a) "no such key: a")
(test/exn (remove-table t-ab 'c) "no such key: c")
\end{lstlisting}
The fifth test is the immutability promise at the Racket level: \code{set-table} made a new hash, and \code{t-ab} still maps \code{a} to 1.
\end{solution}

\subsection{Tables in the Interpreter}
\label{sec:rec6-interp}

\begin{exercise}[The four table cases]
Fill in the four table cases of \code{eval-env}. \code{maketableC} uses \code{make-table}, starting from an empty hash. \code{gettableC}, \code{settableC}, and \code{removetableC} evaluate \code{tbl} and use \code{expect-table} on the result. In \code{settableC}, evaluate \code{tbl} \emph{before} \code{val}, and pass the store from one to the other.
\end{exercise}
\begin{solution}
Each case evaluates its subexpressions while threading the store, then calls a helper:
\begin{lstlisting}
    [maketableC (entries) (make-table env sto entries (hash empty))]
    [gettableC (tbl k)
               (type-case Result (eval-env env sto tbl)
                 [res (v sto-1) (res (get-table (expect-table v) k) sto-1)])]
    [settableC (tbl k val)
               (type-case Result (eval-env env sto tbl)
                 [res (v1 sto-1)
                      (let ([t (expect-table v1)])
                        (type-case Result (eval-env env sto-1 val)
                          [res (v2 sto-2) (res (tableV (set-table t k v2)) sto-2)]))])]
    [removetableC (tbl k)
                  (type-case Result (eval-env env sto tbl)
                    [res (v sto-1) (res (tableV (remove-table (expect-table v) k)) sto-1)])]
\end{lstlisting}
\code{settableC} is the only case with two subexpressions, so it is the only one where order matters. It evaluates \code{tbl} first, then \code{val} in \code{sto-1}. Here \code{expect-table} is called \emph{before} \code{val} is evaluated, so \code{(set-table 5 a (set-box! b 1))} stops before \code{val} has changed the store. Checking after \code{val} would also pass every test below; the difference is only whether the side effects of \code{val} happen before the error.
\end{solution}

\begin{lstlisting}
(test (run '(get-table (make-table (a 1) (b 2)) b)) (numV 2))
(test (run '(get-table (make-table (a (+ 1 2))) a)) (numV 3))
(test (run '(make-table)) (tableV (hash empty)))
(test (run '(make-table (a 1) (b 2))) (tableV t-ab))
(test (run '(get-table (make-table (a 1) (a 2)) a)) (numV 2))   ; the later entry wins
(test/exn (run '(get-table (make-table (a 1)) b)) "no such key: b")

(test (run '(get-table (set-table (make-table (a 1)) b 2) b)) (numV 2))
(test (run '(get-table (set-table (make-table (a 1)) a 2) a)) (numV 2))
(test/exn (run '(get-table (remove-table (make-table (a 1)) a) a)) "no such key: a")
(test/exn (run '(remove-table (make-table (a 1)) b)) "no such key: b")
\end{lstlisting}
Entries are evaluated in the current environment, and can hold any value, here a table and a function:
\begin{lstlisting}
(test (run '(let y 5 (get-table (make-table (a y)) a))) (numV 5))
(test (run '(get-table (get-table (make-table (inner (make-table (x 7)))) inner) x))
      (numV 7))
(test (run '(let t (make-table (inc (lambda x (+ x 1))))
              ((get-table t inc) 41)))
      (numV 42))
\end{lstlisting}
The \code{inc} test shows a table holding a closure, and why \code{parse}'s \code{else} branch must handle a non-symbol head: \code{((get-table t inc) 41)} is an application whose function part is a \code{get-table} expression. \code{make-table} passes the store from one entry to the next, so the second entry sees the first entry's \code{set-box!}; and using something that is not a table as a table is an error:
\begin{lstlisting}
(test (run '(let b (box 0)
              (get-table (make-table (x (set-box! b 1)) (y (unbox b))) y)))
      (numV 1))

(test/exn (run '(get-table 5 a)) "expected a table")
(test/exn (run '(set-table (box 1) a 2)) "expected a table")
(test/exn (run '(remove-table (lambda x x) a)) "expected a table")
\end{lstlisting}

\subsection{Tables vs.\ Boxes}
\label{sec:rec6-vs}

\begin{exercise}[What can change?]
Predict what each program evaluates to before running it.
\begin{lstlisting}
;; a) set-table returns a new table. What happens to t?
(define ex-a
  '(let t (make-table (a 1))
     (begin
       (set-table t a 2)
       (get-table t a))))

;; b) This time we use the table that set-table returns.
(define ex-b
  '(let t (make-table (a 1))
     (get-table (set-table t a 2) a)))

;; c) set! can change what the *variable* t refers to.
(define ex-c
  '(let t (make-table (a 1))
     (begin
       (set! t (set-table t a 2))
       (get-table t a))))

;; d) Two variables, one table. Does u see the change?
(define ex-d
  '(let t (make-table (a 1))
     (let u t
       (begin
         (set! t (set-table t a 2))
         (get-table u a)))))

;; e) A box inside a table.
(define ex-e
  '(let b (box 1)
     (let t (make-table (a b))
       (begin
         (set-box! b 5)
         (unbox (get-table t a))))))

;; f) Which runs first: the table or the value?
(define ex-f
  '(let b (box 0)
     (get-table (set-table (begin (set-box! b 1) (make-table))
                           a
                           (unbox b))
                a)))
\end{lstlisting}
\end{exercise}
\begin{solution}
\begin{lstlisting}
(test (run ex-a) (numV 1))
(test (run ex-b) (numV 2))
(test (run ex-c) (numV 2))
(test (run ex-d) (numV 1))
(test (run ex-e) (numV 5))
(test (run ex-f) (numV 1))
\end{lstlisting}
\begin{center}
\begin{tabular}{@{}lll@{}}
\toprule
ex & result & why \\
\midrule
a & 1 & the new table is thrown away by \code{begin}; \code{t} never changed \\
b & 2 & reads the new table directly \\
c & 2 & \code{set!} puts the new table in \code{t}'s cell \\
d & 1 & \code{u} has its own cell, holding the old table \\
e & 5 & the table holds a \emph{box}; the box's contents changed \\
f & 1 & the table expression runs first, so \code{(unbox b)} sees 1 \\
\bottomrule
\end{tabular}
\end{center}
The way to keep these straight is to ask, for each program, which of the two maps changes.
\begin{lstlisting}[style=repl]
ex-d, just before (get-table u a):

  env:   t -> L0        u -> L1
  store: L0 -> T2       (overridden by set!)
         L1 -> T1       (copied from t when u was bound; set! never touched it)

  T1 = {a: 1}    T2 = {a: 2}      two separate, unchanging tables
\end{lstlisting}
\begin{itemize}
  \item \code{set!} changes the \emph{store}, at the one location its variable names. In \code{ex-c} that is the location we read from; in \code{ex-d} it is not. \code{(let u t ...)} copied the table \emph{value} into \code{u}'s new cell, and that value cannot change.
  \item A table can still \emph{contain} something mutable. In \code{ex-e} the table maps \code{a} to \code{(boxV L)}, and that does not change. What changes is the contents of \code{L} in the store, which is exactly what \code{unbox} reads.
  \item \code{ex-f} is the evaluation-order test of \Cref{sec:rec6-order}, for \code{settableC}. Right to left would give 0.
\end{itemize}
\end{solution}

\begin{exercise}[Tables of locations]
Another design: a table maps each key to a \emph{location} in the store, and the store holds the value. Would that change any of the answers above? When would it matter?
\end{exercise}
\begin{solution}
It depends on what \code{set-table} does with that design.
\begin{itemize}
  \item If \code{set-table} \emph{allocates a fresh location} for \code{k} and copies the other key-to-location pairs, nothing ever writes to an existing table's cell, so all six answers stay the same. The difference is hidden sharing: the old and new tables now point to the \emph{same} cells for every other key. As soon as we add an operation that writes to a key's cell (an in-place \code{update-table!}), a write through one table shows up in the other.
  \item If \code{set-table} \emph{overrides the existing location} for a key that is already present, that is in-place mutation, the case of the next exercise.
\end{itemize}
\end{solution}

\begin{exercise}[Mutable tables]
Which of \code{ex-a} to \code{ex-f} would change if tables could be changed in place, like Racket's mutable hash tables or JavaScript objects?
\end{exercise}
\begin{solution}
\code{ex-a} and \code{ex-d}. In place, \code{set-table} changes the one table that both the old and new names refer to:
\begin{center}
\begin{tabular}{@{}llll@{}}
\toprule
ex & immutable & in place & why it changes \\
\midrule
a & 1 & 2 & the ``thrown away'' result was \code{t} itself \\
d & 1 & 2 & \code{t} and \code{u} alias one table \\
b, c, e, f & same & same & they read the result, or never use \code{set-table} on a shared table \\
\bottomrule
\end{tabular}
\end{center}
\code{ex-d} is the aliasing problem in its usual form. In JavaScript, \code{let u = t; t.a = 2; u.a} gives 2. With immutable tables a copy and a reference behave the same, so aliasing cannot be observed. The table above can be checked by changing \tyname{Table} to a mutable hash (\code{make-hash}, and \code{hash-set!} in \code{make-table} and \code{set-table}).
\end{solution}

\subsection{Challenge: Passing the Store with Monads}

Every multi-subexpression case above has the same shape: evaluate, take the result apart with \code{type-case Result}, pass the new store on, and repeat. That shape can be factored out into two functions, \code{mret} and \code{mbind}. With \tyname{Result} fixed to hold a \tyname{Value}:
\begin{lstlisting}
;; A Stateful is a computation that needs a store to run, and gives back a Result.
(define-type-alias Stateful (Store -> Result))

;; A computation that just returns v, without touching the store.
(define (mret (v : Value)) : Stateful
  (lambda (sto) (res v sto)))

;; Run m, then pass its value to f, and run the computation f returns,
;; using the store m returned.
(define (mbind (m : Stateful) (f : (Value -> Stateful))) : Stateful
  (lambda (sto)
    (type-case Result (m sto)
      [res (v sto-1) ((f v) sto-1)])))

;; eval-env isn't written this way, but a call to it can be turned into a Stateful.
(define (meval (env : Env) (e : Expr)) : Stateful
  (lambda (sto) (eval-env env sto e)))
\end{lstlisting}
A \tyname{Stateful} is a computation that has not run yet: give it a store and it gives back a \tyname{Result}. \code{mbind} is the only place \code{sto-1} appears. It runs \code{m}, hands the value to \code{f}, and runs whatever \code{f} returns \emph{in the store \code{m} left behind}. Threading the store correctly is now \code{mbind}'s job, not each case's.

\begin{exercise}[\code{make-table/m}]
Rewrite \code{make-table} using \code{mret}, \code{mbind}, and \code{meval}: no \code{sto}, and no \code{type-case Result}.
\end{exercise}
\begin{solution}
\begin{lstlisting}
(define (make-table/m (env : Env) (entries : (listof Entry)) (acc : Table)) : Stateful
  (cond
    [(empty? entries) (mret (tableV acc))]
    [else
     (type-case Entry (first entries)
       [entry (k e)
              (mbind (meval env e)
                     (lambda (v)
                       (make-table/m env (rest entries) (hash-set acc k v))))])]))
\end{lstlisting}
Put it next to \code{make-table}: \code{(res (tableV acc) sto)} became \code{(mret (tableV acc))}, and the \code{type-case Result} plus explicit \code{sto-1} became one \code{mbind}. The parameter \code{sto} is gone: \code{make-table/m} returns a computation, and whoever runs it supplies the store.
\end{solution}

\begin{exercise}[\code{set-table/m}]
Rewrite the \code{settableC} case of \code{eval-env} the same way, as \code{(set-table/m env tbl k val)}.
\end{exercise}
\begin{solution}
\begin{lstlisting}
(define (set-table/m (env : Env) (tbl : Expr) (k : symbol) (val : Expr)) : Stateful
  (mbind (meval env tbl)
         (lambda (tv)
           (let ([t (expect-table tv)])
             (mbind (meval env val)
                    (lambda (v) (mret (tableV (set-table t k v)))))))))
\end{lstlisting}
The nesting of the \code{mbind}s \emph{is} the evaluation order: \code{tbl} is outermost, so it runs first. To evaluate right to left, swap the two \code{mbind}s; no store variable would need renaming.
\end{solution}

To test these we need an environment and a store to run them in. \code{env-b} and \code{sto-b} build by hand the state just after \code{(let b (box 0) ...)}: \code{b} lives at $-1$, and that cell holds a box pointing at $-2$, which holds 0. Negative locations cannot clash with the ones \code{new-loc} hands out.
\begin{lstlisting}
(define env-b (extend-env (bind 'b -1) empty-env))
(define sto-b (override-store (cell -1 (boxV -2))
                              (override-store (cell -2 (numV 0)) empty-store)))

;; The second entry sees the first entry's set-box!.
(test (get-table (expect-table
                  (res-v ((make-table/m env-b
                                        (list (entry 'x (parse '(set-box! b 1)))
                                              (entry 'y (parse '(unbox b))))
                                        (hash empty))
                          sto-b)))
                 'y)
      (numV 1))

;; set-table/m evaluates the table before the value, and returns the final store.
(define set-table/m-result
  ((set-table/m env-b
                (parse '(begin (set-box! b 1) (make-table)))
                'a
                (parse '(unbox b)))
   sto-b))
(test (get-table (expect-table (res-v set-table/m-result)) 'a) (numV 1))
(test (fetch -2 (res-s set-table/m-result)) (numV 1))
\end{lstlisting}

\begin{exercise}[Using them in \code{eval-env}]
Use \code{make-table/m} and \code{set-table/m} in the \code{maketableC} and \code{settableC} cases of \code{eval-env}, and check that all the tests still pass.
\end{exercise}
\begin{solution}
A \tyname{Stateful} becomes a \tyname{Result} by applying it to the current store, so the two cases become one line each:
\begin{lstlisting}
    [maketableC (entries) ((make-table/m env entries (hash empty)) sto)]
    [settableC (tbl k val) ((set-table/m env tbl k val) sto)]
\end{lstlisting}
All the table tests pass. One thing to watch: the tests that call \code{run} must come after the \code{define}s of \code{make-table/m} and \code{set-table/m}. Module-level code runs top to bottom, so an earlier test fails with ``cannot reference an identifier before its definition''.
\end{solution}

\subsection{Summary}

\begin{center}
\begin{tabular}{@{}ll@{}}
\toprule
idea & where it shows up \\
\midrule
env: name to location; store: location to value & \code{lookup} then \code{fetch} \\
every binding gets a fresh cell & \code{appC} calls \code{new-loc}; \code{set!} needs it \\
the store is threaded, so order is a choice & \code{ex-order}, \code{ex-f}, the C program \\
tables are values, boxes are references & \code{ex-a} to \code{ex-e} \\
mutable vs.\ immutable is about aliasing & \code{ex-a} and \code{ex-d} flip with in-place tables \\
\code{mret}/\code{mbind} hide the store-passing & \code{make-table/m}, \code{set-table/m} \\
\bottomrule
\end{tabular}
\end{center}
